\subsubsection{Changing the Order of Integration}

\content{
\begin{example} Evaluate the integral 
$$ \int_0^1\int_x^1 \sin(y^2) dydx. $$
\end{example}
}{% presentation
\vspace{4in}
}{% nopresentation
\vspace{2in}
}{% comments
Draw a diagram of the region to switch the bounds to $x$ first. 
}

\content{
\begin{example} Evaluate the integral 
$$ \int_0^1\int_{3y}^3 e^{x^2} \, dxdy. $$
\end{example}
}{% presentation
\vspace{4in}
}{% nopresentation
\vspace{2in}
}{% comments
}

\content{
\begin{theorem} (Properties of Double Integrals) 
\begin{enumerate} 
\item $\iint_D (f_1(x,y) + f_2(x,y))\,dA = \iint_D f_1(x,y)\,dA + \iint_D
f_2(x,y)\, DA. $
\item $\iint_D cf(x,y)\,dA = c\iint_D f(x,y)\,dA. $
\item $\iint_D f(x,y)\, dA = \iint_{D_1} f(x,y)\,dA + \iint_{D_2} f(x,y)\,dA $,
where $D =D_1\cup D_2$ and $D_1$ and $D_2$ do not overlap, except possibly on
their boundaries. 
\item $\iint_D 1\, dA = A(D)$. 
\item If $m\leq f(x,y)\leq M$ on $D$, then 
$mA(D)\leq \iint_D f(x,y)\, dA\leq MA(D). $
\end{enumerate}
\end{theorem}
}{% presentation
}{% nopresentation
}{% comments
}

\subsection{Double Integrals in Polar Coordinates}


\content{

}{% presentation
\vspace{3in}
}{% nopresentation
\vspace{2in}
}{% comments
Review polar coordinates, go over how the extra r shows up in the integration:
\begin{enumerate}
\item draw a basic region
\item draw sectors based on $\Delta \theta$ and $\Delta r$
\item show that the area of each polar rectangle is given by 
$$\Delta A_i = \frac{1}{2}(r_i^2-r_{i-1}^2)\Delta \theta$$
by using the difference of the areas of two sectors of a circle
\item use the fact that $r_i^* = \frac{1}{2}(r_i+r_{i-1})$ is the center point
of each polar rectangle to show that the area is 
$$ \Delta A_i = r_i^* \Delta r \Delta \theta$$
\end{enumerate}
}

\content{
\begin{theorem} If $f$ is a continuous function on a polar rectangle $R$ given
by $a\leq \theta\leq b$ and $\alpha\leq r\leq \beta$, then 
$$ \iint_Rf(x,y)\,dA = \int_\alpha^\beta\int_a^b rf(r\cos\theta,
r\sin\theta)\,d\theta dr. $$
\end{theorem}
}{% presentation
}{% nopresentation
}{% comments
}

\content{
\begin{example} Evaluate the integral 
$$ \iint_R (3x+4y^2)\,dA $$
where $R$ is the region in the upper-half plane bounded by the circles
$x^2+y^2=1$ and $x^2+y^2=4$.
\end{example}
}{% presentation
\vspace{4in}
}{% nopresentation
\vspace{2in}
}{% comments
}

\content{
\begin{example} Find the volume of the solid bounded by the plane $z=0$ and the
surface $z=1-x^2-y^2$.
\end{example}
}{% presentation
\vspace{3in}
}{% nopresentation
\vspace{2in}
\newpage
}{% comments
}

\content{
\begin{example} Use a double integral to find the area of the region enclosed by
one loop of the function $r=\cos 2\theta$.
\end{example}
}{% presentation
\vspace{4in}
}{% nopresentation
\vspace{2in}
}{% comments
}

\content{
\begin{example} FInd the volume of the solid that lies in the cylinder
$x^2+y^2=2x$ and between the $xy$-plane and the surface $z=x^2+y^2$.
\end{example}
}{% presentation
\vspace{4in}
}{% nopresentation
\vspace{2in}
}{% comments
Complete the square in the cylinder equation to find it's center and radius
then set $x=r\cos\theta$ and substitute this into the cylinder equation (and
$r^2=x^2+y^2$)
}
